\section{有限精度・計算量と成立条件}
\label{sec:hhl-assumptions-and-limitations}

HHLアルゴリズムでは，$N$次元の解をおよそ$\log_2N$個の量子ビットで表せる．
また，入力を量子回路へ短時間で与えられる場合には，解に関する特定の期待値を，
$N$に比例しない時間で求められる可能性がある\cite{shimada2020}．
ただし，どの線形方程式でも計算が速くなるわけではない．
HHL全体を評価するには，次の条件を確認する必要がある．

\subsection{行列のエルミート性と可逆性}
\label{subsec:hermiticity-and-invertibility}

本論文で説明したHHLアルゴリズムでは，$A$がエルミートで，
逆行列を持つことを仮定する．
線形回帰の正規方程式で用いる$X^{\mathsf{T}}X$は実対称行列であるため，エルミート性を満たす．
ただし，$X$の列が線形独立でない場合，$X^{\mathsf{T}}X$は逆行列を持たない．
この場合，線形従属な特徴量を取り除く方法や，小さな特異値を無視する方法，
リッジ回帰で正則化項を加える方法などが必要になる．

\subsection{疎性とハミルトニアンシミュレーション}
\label{subsec:sparsity-and-hamiltonian-simulation}

HHLアルゴリズムでは，$e^{\mathrm{i}At}$を短い時間で実行できなければならない．
各行に零でない要素が少ししかない行列を疎行列という．
元のHHLアルゴリズムでは，$A$が疎行列であり，
零でない要素の位置と値を短い時間で読み出せると仮定している\cite{shimada2020}．
しかし，$X$が疎行列であっても，積$X^{\mathsf{T}}X$は密行列となる場合がある．
したがって，訓練データの行列$X$が疎行列であるだけでは，
HHLが短い時間で実行できるとはいえない．実際に使う$A$の形を調べる必要がある．

\subsection{状態準備の計算量}
\label{subsec:state-preparation-cost}

HHLアルゴリズムを始めるには，$\bm{b}$の各成分を振幅に記録した
$\ket{b}$を準備する必要がある．
第\ref{sec:data-encoding}節で述べたように，高次元ベクトルを少数の量子ビットで表せても，
その量子状態を短い時間で作れるとは限らない．
状態準備にベクトルの次元と同程度の時間が必要であれば，HHL部分の利点が相殺される可能性がある．

\subsection{条件数と精度}
\label{subsec:condition-number-and-precision}

条件数$\kappa(A)$が大きい場合，小さな固有値の逆数を高い精度で表現する必要がある．
固有値の近似値を$\widetilde\lambda=\lambda+\Delta$とする．
ここで$\Delta$は近似による誤差であり，
$|\Delta|<|\lambda|$と仮定すると，
\begin{equation}
 \left|\frac1{\widetilde\lambda}-\frac1\lambda\right|
 =\frac{|\Delta|}{|\lambda\widetilde\lambda|}
 \label{eq:reciprocal-eigenvalue-error}
\end{equation}
である．この式から，$|\lambda|$が小さいほど，
同じ大きさの誤差でも逆数の誤差が大きくなることが分かる．
そのため，位相レジスタの量子ビット数を増やし，
行列指数関数と制御回転をより正確に実行する必要がある．
また，補助量子ビットの測定回数も増える．
特に，式\eqref{eq:normal-equation-condition-number}が示すように，
正規方程式では$X$の条件数が二乗される点に注意が必要である．
正則化を行うと条件数を小さくできるが，元の最小二乗問題とは異なる解になる．
計算を安定させることと，元の問題の解に近づけることを両方考える必要がある．

\subsection{出力の読み出し}
\label{subsec:solution-readout}

HHLアルゴリズムの出力$\ket{x}$を測定して，
解ベクトルの全成分を古典データとして得るには，一般に多数の測定が必要になる．
そのため，HHLアルゴリズムは全成分を必要とする用途よりも，
解状態の内積，ある測定結果が得られる確率，または期待値だけを求める用途に適している
\cite{shimada2020}．
最終的に必要な量を少ない測定回数で求められるか，先に確認する必要がある．

\subsection{計算量を評価する範囲}
\label{subsec:end-to-end-evaluation}

計算量は，HHLアルゴリズム単体だけでなく，次の処理を含めて評価する必要がある．

\begin{enumerate}[(1)]
  \item 訓練データから$A$と$\bm{b}$を構成する処理
  \item $A$の要素を読み出し，$e^{\mathrm{i}At}$を実行する処理
  \item $\ket{b}$を準備する処理
  \item HHLアルゴリズムを成功させるための反復
  \item 内積や期待値を必要な精度で求めるための測定
\end{enumerate}

入力の準備から出力の測定までを含む計算量が古典手法より小さい場合に，
HHLを使う利点があるといえる．
したがって，行列の次元だけで比較してはいけない．
行列の零でない要素の数，条件数，状態準備，必要な精度，
出力をどのように使うかを示して比較することが重要である．
